% ============================================================
%  Corriges Chapitre 5 : Mesures des angles et trigonometrie
% ============================================================
\clearpage
\chapter{Mesures des angles et trigonométrie}
\chaptermark{Mesures des angles et trigonométrie}

\begin{corrige}{5.1}
On multiplie par $\dfrac{\pi}{180}$ dans un sens et par $\dfrac{180}{\pi}$ dans l'autre :
\[
210\dgr{} = 210 \times \frac{\pi}{180} = \frac{7\pi}{6}\ \text{rad},
\qquad
240\dgr{} = \frac{4\pi}{3}\ \text{rad},
\qquad
300\dgr{} = \frac{5\pi}{3}\ \text{rad},
\]
et $\dfrac{5\pi}{4}$ rad $= \dfrac{5\pi}{4} \times \dfrac{180}{\pi} = 225\dgr{}$. Contrôles : chaque degré se retrouve en divisant sa fraction par $\dfrac{\pi}{180}$, par exemple $\dfrac{7\pi}{6} \times \dfrac{180}{\pi} = 210$. La longueur d'arc se lit avec $s = r \times \alpha$ :
\[
s = 8 \times \frac{\pi}{4} = 2\pi \approx 6{,}28\ \text{cm}.
\]
\end{corrige}

\begin{corrige}{5.2}
\begin{enumerate}
\item L'échelle est l'hypoténuse d'un triangle rectangle : le mur et le sol en sont les deux côtés de l'angle droit. Pythagore donne $10^{2} = 4^{2} + h^{2}$, d'où $h^{2} = 100 - 16 = 84$ et
\[
h = \sqrt{84} \approx 9{,}2\ \text{m}.
\]
Le sommet de l'échelle touche le mur à environ 9,2 mètres, valeur cohérente : inférieure à la longueur de l'échelle.
\item Le côté de $40$ m est adjacent à l'angle de $25\dgr{}$ et la hauteur $h$ lui est opposée :
\[
\tan 25\dgr{} = \frac{h}{40}
\quad\Longrightarrow\quad
h = 40 \times \tan 25\dgr{} \approx 40 \times 0{,}466 \approx 18{,}7\ \text{m}.
\]
La tour culmine à environ 18,7 mètres. Dans les deux questions, l'hypoténuse (l'échelle) ou l'angle (la visée) est connu : on choisit Pythagore quand les trois côtés sont concernés, et un rapport trigonométrique quand un angle et deux côtés sont liés.
\end{enumerate}
\end{corrige}

\begin{corrige}{5.3}
\begin{enumerate}
\item L'angle $\dfrac{5\pi}{6}$ rad vaut $150\dgr{}$ : son point image se trouve dans le deuxième quadrant, où l'abscisse est négative et l'ordonnée positive. On lit
\[
\cos\frac{5\pi}{6} = -\frac{\sqrt{3}}{2}
\qquad\text{et}\qquad
\sin\frac{5\pi}{6} = \frac{1}{2} .
\]
\item La parité du cosinus (fonction paire) et du sinus (fonction impaire) donne
\[
\cos\left(-\frac{\pi}{3}\right) = \cos\frac{\pi}{3} = \frac{1}{2}
\qquad\text{et}\qquad
\sin\left(-\frac{\pi}{3}\right) = -\sin\frac{\pi}{3} = -\frac{\sqrt{3}}{2} .
\]
\item La périodicité de $2\pi$ permet de retirer un tour complet : $\dfrac{7\pi}{3} = \dfrac{7\pi}{3} - 2\pi = \dfrac{\pi}{3}$, donc
\[
\sin\frac{7\pi}{3} = \sin\frac{\pi}{3} = \frac{\sqrt{3}}{2} .
\]
\end{enumerate}
\end{corrige}

\begin{corrige}{5.4}
\begin{enumerate}
\item Pour $\cos x = \dfrac{\sqrt{2}}{2}$, l'angle du tableau des valeurs remarquables est $\alpha = \dfrac{\pi}{4}$ ; les solutions sont le point image de $\alpha$ et son symétrique par rapport à l'axe des abscisses :
\[
x = \frac{\pi}{4} + 2k\pi \qquad\text{ou}\qquad x = -\frac{\pi}{4} + 2k\pi, \qquad k \in \Z .
\]
\item Pour $\sin x = -\dfrac{1}{2}$, l'angle de référence est $\alpha = -\dfrac{\pi}{6}$ (le sinus est négatif dans les troisième et quatrième quadrants) ; les solutions sont $x = \alpha + 2k\pi$ et son symétrique $x = (\pi - \alpha) + 2k\pi$ :
\[
x = -\frac{\pi}{6} + 2k\pi \qquad\text{ou}\qquad x = \frac{7\pi}{6} + 2k\pi, \qquad k \in \Z .
\]
\item Pour $\cos x = 1$, l'angle solution dans $[0;\ 2\pi[$ est $\alpha = 0$, et son symétrique $-\alpha = 0$ coïncide avec lui : il reste une seule famille de solutions, un point par tour complet :
\[
x = 2k\pi, \qquad k \in \Z .
\]
Dans chaque cas, la périodicité ajoute un tour complet par valeur de $k$ : les équations trigonométriques ont une infinité de solutions dans $\R$.
\end{enumerate}
\end{corrige}
